\documentclass[11pt,a4paper]{article}

% ---- Packages ----
\usepackage[utf8]{inputenc}
\usepackage[T1]{fontenc}
\usepackage{amsmath,amssymb,amsthm}
\usepackage{mathtools}
\usepackage{algorithm}
\usepackage{algpseudocode}
\usepackage{booktabs}
\usepackage{multirow}
\usepackage{graphicx}
\usepackage{xcolor}
\usepackage{hyperref}
\usepackage{cleveref}
\usepackage[margin=2.5cm]{geometry}
\usepackage{enumitem}
\usepackage{tikz}
\usetikzlibrary{matrix,positioning,arrows.meta,fit,backgrounds}
\usepackage{caption}
\usepackage{subcaption}

% ---- Macros ----
\newcommand{\setA}{\mathcal{A}}
\newcommand{\setU}{\mathcal{U}}
\newcommand{\setS}{\mathcal{S}}
\newcommand{\setT}{\mathcal{T}}
\newcommand{\setQ}{\mathcal{Q}}
\newcommand{\SU}{\mathrm{SU}}
\newcommand{\SEQ}{\mathrm{SEQ}}
\newcommand{\TSA}{\mathrm{TSA}}
\newcommand{\TSS}{\mathrm{TSS}}
\newcommand{\HN}{\mathrm{HN}}
\newcommand{\HE}{\mathrm{HE}}

\newtheorem{definition}{Definition}
\newtheorem{example}{Example}

\title{\textbf{Home Care Optimised Resource Allocation Problem (HCORAP)}\\[6pt]
\large Problem Description, Mathematical Models, and Dataset}
\author{}
\date{}

\begin{document}
\maketitle

\tableofcontents
\newpage

%=============================================================================
\section{Problem Description}\label{sec:problem}
%=============================================================================

The \textbf{Home Care Optimised Resource Allocation Problem (HCORAP)} arises in
the context of home care service companies that assign care workers (agents) to
perform services for users (patients/clients) over a weekly planning horizon.

A home care company manages a set of \emph{agents} $\setA = \{a_1, \ldots, a_A\}$
who visit a set of \emph{users} $\setU = \{u_1, \ldots, u_U\}$.
Each user requires a set of \emph{services}
$\setS = \{s_1, \ldots, s_S\}$ to be performed during the week,
which is discretised into a set of \emph{time slots}
$\setT = \{t_1, \ldots, t_{T}\}$ (e.g.\ $T = 60$ for 12~slots per day
$\times$ 5~working days).

The problem asks to find an assignment of agents to services at specific time slots
that satisfies a collection of hard constraints while maximising a multi-criteria
objective function composed of:
\begin{enumerate}[label=(\roman*)]
    \item \textbf{Similarity (revenue):} the accumulated expertise reward of
          assigning agents to the services they are best suited for;
    \item \textbf{Continuity of care:} a preference for having as few distinct
          agents as possible attending the services of a single user that share
          the same service type;
    \item \textbf{Working-hour cost:} a penalty incurred whenever an agent works
          beyond their contractual (normal) hours.
\end{enumerate}

\subsection{Input Data}

An HCORAP instance is defined by the following data.

\begin{itemize}
    \item $U$: number of users; \;
          $S$: number of services; \;
          $A$: number of agents; \;
          $T$: number of time slots.
    \item $\SU_u \subseteq \setS$: the set of services belonging to user $u$.
          These form a partition of $\setS$.
    \item $\SEQ_q \subseteq \setS$ ($q = 1,\ldots,Q$): groups of services that
          belong to the same user \emph{and} have the same service type.
          Ideally, all services in a $\SEQ_q$ should be assigned to the same agent
          (continuity of care).
    \item $\TSA_{a,t} \in \{0,1\}$: availability matrix;
          $\TSA_{a,t} = 1$ iff agent $a$ is available to work at time slot $t$.
    \item $\TSS_{s,t} \in \{0,1\}$: suitability matrix;
          $\TSS_{s,t} = 1$ iff service $s$ can be performed at time slot $t$.
    \item $r_{a,s} \in \{0,1,2,3,4\}$: reward (expertise) score of assigning agent
          $a$ to service $s$.  A value of $0$ means the assignment is
          \emph{forbidden} (the agent lacks the required qualification).
    \item $\HN_a \in \mathbb{Z}_{\ge 0}$: contractual (normal) weekly working
          hours for agent $a$.
    \item $\HE_a \in \mathbb{Z}_{\ge 0}$: maximum allowed extra working hours
          for agent $a$ (beyond $\HN_a$).
    \item $P \in \mathbb{Z}$: penalty weight for each extra working hour
          (typically $P = -1$).
\end{itemize}


%=============================================================================
\section{Illustrative Example}\label{sec:example}
%=============================================================================

\begin{example}
Consider a small instance with $U=2$ users, $A=2$ agents,
$S=4$ services, and $T=4$ time slots per week (for illustration purposes).

\medskip\noindent
\textbf{Service--user assignment:}
\[
    \SU_1 = \{s_1, s_2\}, \quad \SU_2 = \{s_3, s_4\}.
\]

\noindent
\textbf{Continuity-of-care groups:}
\[
    \SEQ_1 = \{s_1, s_2\}, \quad \SEQ_2 = \{s_3\}, \quad \SEQ_3 = \{s_4\}.
\]
Here $s_1$ and $s_2$ are services of the same type for user~1, so they form
one sequence.

\medskip\noindent
\textbf{Agent availability} ($\TSA$) \textbf{and service time windows} ($\TSS$):

\[
\TSA = \begin{pmatrix}
1 & 1 & 0 & 0 \\
0 & 1 & 1 & 1
\end{pmatrix}, \quad
\TSS = \begin{pmatrix}
1 & 1 & 0 & 0 \\
0 & 1 & 1 & 0 \\
1 & 0 & 1 & 0 \\
0 & 0 & 1 & 1
\end{pmatrix}.
\]
Agent~$a_1$ is available in slots $t_1,t_2$; agent~$a_2$ in slots $t_2,t_3,t_4$.
Service~$s_1$ can be done in $t_1,t_2$; service $s_2$ in $t_2,t_3$; etc.

\medskip\noindent
\textbf{Reward matrix} ($r$):
\[
r = \begin{pmatrix}
3 & 2 & 0 & 1 \\
1 & 0 & 4 & 3
\end{pmatrix}.
\]
Agent~$a_1$ is forbidden from service~$s_3$ ($r_{1,3}=0$).
Agent~$a_2$ is forbidden from service~$s_2$ ($r_{2,2}=0$).

\medskip\noindent
\textbf{Working hours:}
$\HN_1 = 2$, $\HE_1 = 0$; $\HN_2 = 2$, $\HE_2 = 1$; $P = -1$.

\medskip\noindent
\textbf{A feasible solution:}
\begin{center}
\begin{tabular}{cccc}
\toprule
Service & Agent & Time slot & Reward \\
\midrule
$s_1$ & $a_1$ & $t_1$ & 3 \\
$s_2$ & $a_1$ & $t_2$ & 2 \\
$s_3$ & $a_2$ & $t_3$ & 4 \\
$s_4$ & $a_2$ & $t_4$ & 3 \\
\bottomrule
\end{tabular}
\end{center}

\textbf{Analysis:}
\begin{itemize}
    \item \emph{Similarity (revenue):} $3 + 2 + 4 + 3 = 12$.
    \item \emph{Continuity of care:} $\SEQ_1 = \{s_1,s_2\}$ is served entirely by
          $a_1$ (1~agent $\to$ penalty~$0$). Total continuity penalty $= 0$.
    \item \emph{Working-hour cost:} agent~$a_1$ works 2~services
          $= \HN_1 \Rightarrow$ no extra hours.
          Agent~$a_2$ works 2~services $= \HN_2 \Rightarrow$ no extra hours.
    \item \textbf{Total objective}: $12 - 0 - 0 = 12$ (to be maximised).
\end{itemize}
\end{example}


%=============================================================================
\section{Main Constraints}\label{sec:constraints}
%=============================================================================

The HCORAP imposes the following hard constraints:

\begin{enumerate}[label=\textbf{C\arabic*.}]
    \item \label{c:amo-service}
    \textbf{At most one assignment per service.}
    Each service is attended by at most one agent at at most one time slot.

    \item \label{c:amo-agent-ts}
    \textbf{At most one service per agent per time slot.}
    An agent can perform at most one service in any given time slot.

    \item \label{c:agent-availability}
    \textbf{Agent availability.}
    Agent $a$ may only work at time slot $t$ if $\TSA_{a,t} = 1$.

    \item \label{c:amo-user-ts}
    \textbf{At most one service per user per time slot.}
    A user cannot receive two services simultaneously.

    \item \label{c:service-coverage}
    \textbf{Service coverage.}
    Every service must be attended at some time slot (hard constraint in the
    basic model; can be relaxed to a soft constraint in the stratified variant).

    \item \label{c:max-hours}
    \textbf{Maximum working hours.}
    Agent $a$ may work at most $\HN_a + \HE_a$ services in total.

    \item \label{c:qualification}
    \textbf{Qualification.}
    Agent $a$ can only be assigned to service $s$ if $r_{a,s} > 0$.
\end{enumerate}

\noindent
Additionally, the following \textbf{soft (optimisation) objectives} are encoded:

\begin{enumerate}[label=\textbf{O\arabic*.}]
    \item \textbf{Maximise total expertise reward:}
    $\displaystyle\sum_{a \in \setA}\sum_{s \in \setS} r_{a,s} \cdot y_{a,s}$.

    \item \textbf{Minimise continuity-of-care violation:}
    For each sequence $\SEQ_q$, minimise the number of distinct agents assigned
    to its services (ideally 1).

    \item \textbf{Minimise extra-hour penalty:}
    For each agent $a$, penalise by $|P|$ for every service beyond $\HN_a$.
\end{enumerate}


%=============================================================================
\section{Integer Linear Programming (ILP) Model}\label{sec:ilp}
%=============================================================================

\subsection{Decision Variables}

\begin{align}
    x_{a,s,t} &\in \{0,1\}
    && \text{$= 1$ iff agent $a$ performs service $s$ at time slot $t$} \label{eq:x}\\
    y_{a,s} &\in \{0,1\}
    && \text{$= 1$ iff agent $a$ performs service $s$ (at any time slot)} \label{eq:y}\\
    z_{s,t} &\in \{0,1\}
    && \text{$= 1$ iff service $s$ is performed at time slot $t$} \label{eq:z}\\
    w_{a,q} &\in \{0,1\}
    && \text{$= 1$ iff agent $a$ performs some service in $\SEQ_q$} \label{eq:w}\\
    e_a &\in \mathbb{Z}_{\ge 0}
    && \text{number of extra hours worked by agent $a$} \label{eq:e}
\end{align}

\subsection{Objective Function}

\begin{equation}\label{eq:obj}
\boxed{%
\max \;\;
\underbrace{\sum_{a \in \setA}\sum_{s \in \setS} r_{a,s} \cdot y_{a,s}}_{\text{expertise reward}}
\;-\;
\underbrace{\sum_{q=1}^{Q} \biggl(\sum_{a \in \setA} w_{a,q} - 1\biggr)^{+}}_{\text{continuity penalty}}
\;+\;
\underbrace{P \cdot \sum_{a \in \setA} e_a}_{\text{extra-hour penalty (}P < 0\text{)}}
}
\end{equation}

\noindent
where $(z)^{+} = \max(z, 0)$.  In practice, the continuity penalty is linearised
using auxiliary variables (see the MaxSAT encoding below for a sorting-network
approach), or equivalently: for each $\SEQ_q$, minimise
$\sum_{a \in \setA} w_{a,q}$.

\subsection{Constraints}

\paragraph{Linking $x$ and $y$:}
\begin{align}
    y_{a,s} &= \sum_{t \in \setT} x_{a,s,t}
    && \forall a \in \setA,\; s \in \setS \label{eq:link-xy}
\end{align}

\paragraph{Linking $x$ and $z$:}
\begin{align}
    z_{s,t} &= \sum_{a \in \setA} x_{a,s,t}
    && \forall s \in \setS,\; t \in \setT \label{eq:link-xz}
\end{align}

\paragraph{C1 -- At most one assignment per service:}
\begin{align}
    \sum_{a \in \setA}\sum_{t \in \setT} x_{a,s,t} &\le 1
    && \forall s \in \setS \label{eq:c1}
\end{align}

\paragraph{C2 -- At most one service per agent per time slot:}
\begin{align}
    \sum_{s \in \setS} x_{a,s,t} &\le 1
    && \forall a \in \setA,\; t \in \setT \label{eq:c2}
\end{align}

\paragraph{C3 -- Agent availability and service suitability:}
\begin{align}
    x_{a,s,t} &= 0
    && \forall (a,s,t) \text{ s.t. } \TSA_{a,t} = 0 \text{ or }
       \TSS_{s,t} = 0 \text{ or } r_{a,s} = 0 \label{eq:c3}
\end{align}

\paragraph{C4 -- At most one service per user per time slot:}
\begin{align}
    \sum_{s \in \SU_u} z_{s,t} &\le 1
    && \forall u \in \setU,\; t \in \setT \label{eq:c4}
\end{align}

\paragraph{C5 -- Service coverage (hard):}
\begin{align}
    \sum_{t \in \setT} z_{s,t} &\ge 1
    && \forall s \in \setS \label{eq:c5}
\end{align}

\paragraph{C6 -- Maximum working hours:}
\begin{align}
    \sum_{s \in \setS} y_{a,s} &\le \HN_a + \HE_a
    && \forall a \in \setA \label{eq:c6-hard}
\end{align}

\paragraph{Extra hours:}
\begin{align}
    e_a &\ge \sum_{s \in \setS} y_{a,s} - \HN_a
    && \forall a \in \setA \label{eq:extra}\\
    e_a &\ge 0
    && \forall a \in \setA  \notag
\end{align}

\paragraph{Linking $y$ and $w$:}
\begin{align}
    w_{a,q} &\ge y_{a,s}
    && \forall a \in \setA,\; q = 1,\ldots,Q,\; s \in \SEQ_q \label{eq:link-yw1}\\
    w_{a,q} &\le \sum_{s \in \SEQ_q} y_{a,s}
    && \forall a \in \setA,\; q = 1,\ldots,Q \label{eq:link-yw2}
\end{align}


%=============================================================================
\section{SAT / MaxSAT Model}\label{sec:sat}
%=============================================================================

The paper~\cite{collcaballero2025hcorap} proposes a \textbf{partial weighted
MaxSAT} encoding.  All hard constraints are encoded as \emph{hard clauses},
while the three optimisation objectives are encoded as \emph{soft clauses} with
appropriate weights.

\subsection{Boolean Variables}

The following Boolean variables are created (only for \emph{valid} combinations,
i.e.\ where $r_{a,s} > 0$, $\TSA_{a,t} = 1$, and $\TSS_{s,t} = 1$):

\begin{itemize}
    \item $x_{a,s,t}$: true iff agent $a$ performs service $s$ at time slot $t$.
    \item $y_{a,s}$: true iff agent $a$ performs service $s$ at some time slot.
    \item $z_{s,t}$ (\texttt{su} in code): true iff service $s$ is
          performed at time slot $t$.
    \item $w_{a,q}$ (\texttt{s} in code): true iff agent $a$ performs at
          least one service in $\SEQ_q$.
\end{itemize}

\subsection{Hard Clauses}

\paragraph{Reification (channelling):}
\begin{align}
    y_{a,s} &\Leftrightarrow \bigvee_{t \in \setT} x_{a,s,t}
    && \text{(encoded via unit implication + one big clause)} \label{eq:reif-y}\\
    z_{s,t} &\Leftrightarrow \bigvee_{a \in \setA} x_{a,s,t}
    && \text{(similarly)} \label{eq:reif-z}\\
    w_{a,q} &\Leftrightarrow \bigvee_{s \in \SEQ_q} y_{a,s}
    && \text{(for $|\SEQ_q| > 1$; otherwise $w_{a,q} = y_{a,s}$)}
    \label{eq:reif-w}
\end{align}

\paragraph{At-Most-One (AMO) constraints:}
These are encoded using standard AMO encodings (e.g.\ ladder/commander encoding):
\begin{align}
    &\mathrm{AMO}\bigl(\{x_{a,s,t} : a \in \setA,\, t \in \setT\}\bigr)
    && \forall s \in \setS \quad \text{(C1)} \label{eq:amo1}\\
    &\mathrm{AMO}\bigl(\{x_{a,s,t} : s \in \setS\}\bigr)
    && \forall a \in \setA,\; t \in \setT \quad \text{(C2)} \label{eq:amo2}\\
    &\mathrm{AMO}\bigl(\{z_{s,t} : s \in \SU_u\}\bigr)
    && \forall u \in \setU,\; t \in \setT \quad \text{(C4)} \label{eq:amo4}
\end{align}

\paragraph{Implicit constraint C3:}
Handled during variable creation: if
$r_{a,s} = 0$, $\TSA_{a,t} = 0$, or $\TSS_{s,t} = 0$,
the corresponding $x_{a,s,t}$ variable is set to a constant ``false'' literal.

\paragraph{Service coverage (C5, hard version):}
\begin{equation}\label{eq:cov-hard}
    \bigvee_{t \in \setT} z_{s,t} \qquad \forall s \in \setS
\end{equation}

\paragraph{Maximum working hours (C6):}
Encoded using \textbf{sorting networks}.  For each agent $a$, the vector of
$y_{a,s}$ variables is sorted; then the output at position
$\HN_a + \HE_a$ is forced to false:
\begin{equation}\label{eq:max-hours}
    \mathrm{Sort}\bigl(\{y_{a,s} : s \in \setS,\, r_{a,s} > 0\}\bigr)
    \;\Longrightarrow\;\;
    \neg\, \mathrm{out}[\HN_a + \HE_a]
    \qquad \forall a \in \setA
\end{equation}

\subsection{Soft Clauses}

\paragraph{O1 -- Expertise reward:}
For each valid pair $(a,s)$ with $r_{a,s} > 0$:
\begin{equation}\label{eq:soft-reward}
    \bigl(y_{a,s},\;\; \text{weight} = r_{a,s}\bigr) \qquad \text{soft clause}
\end{equation}
The MaxSAT solver maximises the total weight of satisfied soft clauses.

\paragraph{O2 -- Continuity of care:}
For each sequence $\SEQ_q$ with $|\SEQ_q| > 1$, a sorting network is applied
to the vector $(w_{1,q}, w_{2,q}, \ldots, w_{A,q})$.
The sorted outputs $\mathrm{out}[0], \ldots, \mathrm{out}[p-1]$
(where $p = \min(A, |\SEQ_q|)$)
are penalised with unit-weight soft clauses:
\begin{equation}\label{eq:soft-cont}
    \bigl(\neg\,\mathrm{out}[i],\;\; \text{weight} = 1\bigr)
    \quad i = 0,\ldots,p-1, \qquad \text{for each } q
\end{equation}
This effectively penalises each additional agent working in the same sequence.

\paragraph{O3 -- Extra-hour penalty:}
For each agent $a$ whose potential workload exceeds $\HN_a$,
the sorted output positions from $\HN_a$ to
$\min(\HN_a + \HE_a, |\{s : r_{a,s} > 0\}|) - 1$
are penalised:
\begin{equation}\label{eq:soft-extra}
    \bigl(\neg\,\mathrm{out}[k],\;\; \text{weight} = |P|\bigr)
    \quad k = \HN_a,\ldots \qquad \text{for each } a
\end{equation}
Note that the penalty weight $|P|$ is negative in the objective (since $P < 0$),
meaning it is a cost to be subtracted from the total reward.

\subsection{Stratified Variant}

In the \textbf{stratified} encoding, constraint~C5 (service coverage) becomes
a \emph{soft} clause with a weight high enough to be strictly prioritised over
the other objectives:
\begin{equation}\label{eq:cov-soft}
    \biggl(\bigvee_{t \in \setT} z_{s,t},\;\;
    \text{weight} = W_{\text{total}} + 1\biggr)
    \qquad \forall s \in \setS
\end{equation}
where $W_{\text{total}}$ is the sum of all other soft-clause weights.
This creates a lexicographic optimisation: first maximise the number of covered
services, then maximise the original objective.


%=============================================================================
\section{Optimisation Approaches}\label{sec:optimisation}
%=============================================================================

The HCORAP is a \textbf{Constrained Optimisation Problem (COP)}: find an
assignment satisfying all hard constraints (Section~\ref{sec:constraints}) that
\emph{maximises} the weighted objective
(Equation~\ref{eq:obj}).  Since standard SAT solvers handle only
satisfiability, optimisation must be achieved by a meta-algorithm that
orchestrates one or more SAT calls.  In what follows we present three
paradigms---\emph{Non-Incremental SAT}, \emph{Incremental SAT}, and
\emph{MaxSAT}---and compare their trade-offs in the context of HCORAP.

We write $\varphi$ for the conjunction of all hard clauses
(Section~\ref{sec:sat}), $W$ for the total soft-clause weight budget,
and $\mathrm{obj}(\sigma)$ for the objective value of a satisfying
assignment $\sigma$.
A \emph{bounded} SAT query $\mathrm{SAT}(\varphi, \mathrm{obj} \ge K)$
asks whether there exists a model of $\varphi$ whose objective is at
least $K$.  This is realised by augmenting $\varphi$ with a
\textbf{cardinality / pseudo-Boolean constraint} that enforces the bound.

%-----------------------------------------------------------------------------
\subsection{Non-Incremental SAT}\label{sec:opt-nonincr}
%-----------------------------------------------------------------------------

In the non-incremental approach, a \textbf{fresh SAT solver instance} is
created for every query.  The formula $\varphi$ and the cardinality encoding
for the current bound are regenerated from scratch each time.

\subsubsection{Top-Down Linear Search}

Start from any feasible solution and iteratively tighten the lower bound until
the problem becomes unsatisfiable.

\begin{algorithm}[H]
\caption{Non-Incremental Top-Down Linear Search (maximisation)}
\label{alg:topdown}
\begin{algorithmic}[1]
\State $\sigma^* \leftarrow \textsc{SAT}(\varphi)$
    \Comment{find any feasible solution}
\State $K \leftarrow \mathrm{obj}(\sigma^*)$
\While{$\textsc{SAT}(\varphi \,\wedge\, \mathrm{obj} \ge K\!+\!1)$}
    \Comment{fresh solver each call}
    \State $\sigma^* \leftarrow$ model;
           \;$K \leftarrow \mathrm{obj}(\sigma^*)$
\EndWhile
\State \Return $\sigma^*$     \Comment{optimal}
\end{algorithmic}
\end{algorithm}

\textbf{Number of SAT calls:} $\mathrm{obj}^* - \mathrm{obj}_0 + 2$,
where $\mathrm{obj}_0$ is the first solution's value.
In the worst case this can be $O(W)$, making it impractical for large $W$.

\subsubsection{Bottom-Up Linear Search}

Start from $K = W$ (maximum possible objective) and decrease until satisfiable.

\begin{algorithm}[H]
\caption{Non-Incremental Bottom-Up Linear Search (maximisation)}
\label{alg:bottomup}
\begin{algorithmic}[1]
\State $K \leftarrow W$
\While{$\neg\,\textsc{SAT}(\varphi \,\wedge\, \mathrm{obj} \ge K)$}
    \Comment{fresh solver}
    \State $K \leftarrow K - 1$
\EndWhile
\State \Return model       \Comment{optimal with $\mathrm{obj} = K$}
\end{algorithmic}
\end{algorithm}

\textbf{Number of SAT calls:} $W - \mathrm{obj}^* + 1$.
Efficient when $\mathrm{obj}^*$ is close to $W$; poor otherwise.

\subsubsection{Binary (Dichotomic) Search}

The most efficient non-incremental strategy.  Maintain an interval $[lb, ub]$
and probe the midpoint.

\begin{algorithm}[H]
\caption{Non-Incremental Binary Search (maximisation)}
\label{alg:binsearch}
\begin{algorithmic}[1]
\State $lb \leftarrow 0$;\;\; $ub \leftarrow W$;\;\; $\sigma^* \leftarrow \text{None}$
\While{$lb \le ub$}
    \State $\mathit{mid} \leftarrow \lceil(lb + ub)/2\rceil$
    \If{$\textsc{SAT}(\varphi \,\wedge\, \mathrm{obj} \ge \mathit{mid})$}
        \Comment{fresh solver}
        \State $\sigma^* \leftarrow$ model;\;\;
               $lb \leftarrow \mathrm{obj}(\sigma^*)$
        \Comment{tighten using actual obj.\ value}
    \Else
        \State $ub \leftarrow \mathit{mid} - 1$
    \EndIf
\EndWhile
\State \Return $\sigma^*$     \Comment{optimal}
\end{algorithmic}
\end{algorithm}

\textbf{Number of SAT calls:}
$O\!\bigl(\log_2(W)\bigr)$.
For a typical HCORAP instance with $W \approx 600$, this gives
$\lceil \log_2 600 \rceil = 10$ calls.

\subsubsection{Cardinality Encoding}

The bound $\mathrm{obj} \ge K$ is enforced by encoding the (weighted) sum of
soft-clause satisfaction as a cardinality or pseudo-Boolean constraint.
Common encodings include~\cite{sinz2005cardinality,bailleux2003totalizer}:

\begin{itemize}
    \item \textbf{Sequential counter}~\cite{sinz2005cardinality}:
          $O(n \cdot K)$ auxiliary variables and clauses.  Compact with good
          unit-propagation properties.
    \item \textbf{Totalizer}~\cite{bailleux2003totalizer}:
          $O(n \log n)$ variables, $O(n \log^2 n)$ clauses.
          Better for large $n$.
    \item \textbf{Sorting networks}~\cite{asin2011cardinality}: based on
          odd--even merge, balances encoding size and propagation strength.
          \emph{This is the encoding used in the HCORAP implementation.}
\end{itemize}

In the non-incremental setting, the cardinality encoding is
\textbf{rebuilt from scratch} for each SAT call, which is the main
source of overhead.

\subsubsection{Key Properties}

\begin{center}
\begin{tabular}{lc}
\toprule
\textbf{Property} & \textbf{Non-Incremental SAT} \\
\midrule
Learned clauses reused across iterations & No \\
Cardinality encoding rebuilt each call   & Yes \\
Best search strategy                     & Binary search \\
Number of SAT calls (binary)             & $O(\log W)$ \\
Implementation complexity                & Low \\
\bottomrule
\end{tabular}
\end{center}


%-----------------------------------------------------------------------------
\subsection{Incremental SAT}\label{sec:opt-incr}
%-----------------------------------------------------------------------------

In the incremental approach, a \textbf{single solver instance} is kept alive
across all iterations.  The solver retains its \emph{learned-clause database},
which can dramatically speed up subsequent queries since conflict analysis from
earlier iterations remains valid.

The key technical challenge is modifying the objective bound between calls
without invalidating the solver state.  Two mechanisms are available:

\subsubsection{Assumption-Based Bound Tightening}

Modern SAT solvers support the \textbf{IPASIR} interface~\cite{ipasir}, which
allows passing \emph{assumption literals} to each \textsc{Solve} call.
Assumptions are temporary unit-clause assertions that are active only for the
current call; they do not permanently alter the clause database.

\paragraph{Idea.}
Encode the sorting network / cardinality constraint \textbf{once} at
initialisation.  Let $o_0, o_1, \ldots, o_n$ be the sorted output variables,
where $o_k = \mathit{true}$ means ``at least $k{+}1$ soft clauses are
satisfied''.  To enforce $\mathrm{obj} \ge K$, simply pass assumption
$o_{K-1} = \mathit{true}$ (or equivalently, for ``$\mathrm{obj} \le K$'',
assume $\neg\, o_K$).

\begin{algorithm}[H]
\caption{Incremental Binary Search with Assumptions (maximisation)}
\label{alg:incr-binary}
\begin{algorithmic}[1]
\State $\mathit{solver} \leftarrow \textsc{NewSolver}(\varphi)$
\State Build sorting network once;
       let $o_0, \ldots, o_n$ be output variables
\State $lb \leftarrow 0$;\;\; $ub \leftarrow W$;\;\; $\sigma^* \leftarrow \text{None}$
\While{$lb \le ub$}
    \State $\mathit{mid} \leftarrow \lceil(lb + ub)/2\rceil$
    \State $\mathit{assumptions} \leftarrow \{o_{\mathit{mid}-1}\}$
        \Comment{enforce obj $\ge$ mid}
    \If{$\mathit{solver}.\textsc{Solve}(\mathit{assumptions})$}
        \State $\sigma^* \leftarrow$ model;\;\;
               $lb \leftarrow \mathrm{obj}(\sigma^*)$
    \Else
        \State $ub \leftarrow \mathit{mid} - 1$
    \EndIf
\EndWhile
\State \Return $\sigma^*$
\end{algorithmic}
\end{algorithm}

\textbf{Advantages:}
\begin{itemize}
    \item The cardinality encoding is built \textbf{once}, not rebuilt each call.
    \item All learned clauses from previous iterations remain in the solver.
    \item Assumptions add zero permanent clauses---the solver's state stays clean.
    \item Typically \textbf{2--10$\times$ faster} than non-incremental binary
          search on the same instance.
\end{itemize}

\subsubsection{Clause-Addition Linear Search}

An alternative incremental approach: after finding a solution with
$\mathrm{obj} = K$, add a \emph{permanent} clause
$(\mathrm{obj} \ge K\!+\!1)$ to the solver and re-solve.

\begin{algorithm}[H]
\caption{Incremental Linear Search with Clause Addition (maximisation)}
\label{alg:incr-linear}
\begin{algorithmic}[1]
\State $\mathit{solver} \leftarrow \textsc{NewSolver}(\varphi)$
\While{$\mathit{solver}.\textsc{Solve}()$}
    \State $\sigma^* \leftarrow$ model;\;\;
           $K \leftarrow \mathrm{obj}(\sigma^*)$
    \State $\mathit{solver}.\textsc{AddClauses}\bigl(\mathrm{Encode}(\mathrm{obj} \ge K\!+\!1)\bigr)$
\EndWhile
\State \Return $\sigma^*$     \Comment{last SAT model is optimal}
\end{algorithmic}
\end{algorithm}

\textbf{Caveat:} Added clauses are \emph{permanent} and cannot be retracted.
Old (subsumed) cardinality constraints accumulate in the clause database,
increasing memory usage.  The assumption-based approach
(Algorithm~\ref{alg:incr-binary}) avoids this issue entirely.

\subsubsection{Application to HCORAP}

The HCORAP codebase implements incremental solving through the
\texttt{DicoOptimizer} class, which supports both modes:
\begin{itemize}
    \item When \texttt{useAssumptions = false}: a fresh encoding is created per
          call (non-incremental binary search, Section~\ref{sec:opt-nonincr}).
    \item When \texttt{useAssumptions = true}: the solver is initialised once
          via \texttt{initAssumptionOptimization}, and subsequent calls use
          \texttt{checkSATAssuming} (incremental binary search,
          Algorithm~\ref{alg:incr-binary}).
\end{itemize}
The sorting-network outputs serve as the assumption handles: assuming a sorted
output to be true or false directly controls the objective bound.

\subsubsection{Key Properties}

\begin{center}
\begin{tabular}{lcc}
\toprule
\textbf{Property} & \textbf{Assumption-based} & \textbf{Clause-addition} \\
\midrule
Learned clauses reused     & Yes  & Yes \\
Encoding built once        & Yes  & No (rebuilt each tightening)\\
Subsumed clause accumulation & None & Yes (growing database) \\
Best search strategy       & Binary search & Linear search \\
\bottomrule
\end{tabular}
\end{center}


%-----------------------------------------------------------------------------
\subsection{MaxSAT}\label{sec:opt-maxsat}
%-----------------------------------------------------------------------------

The most natural approach for HCORAP: encode the problem directly as a
\textbf{partial weighted MaxSAT} instance and delegate the entire optimisation
to a specialised MaxSAT solver.

\subsubsection{Formulation}

A \emph{partial weighted MaxSAT} (PWMS) instance consists of:
\begin{itemize}
    \item \textbf{Hard clauses} $\varphi_H$: must be satisfied (weight $= \infty$).
    \item \textbf{Soft clauses} $\{(c_i, w_i)\}_{i=1}^{m}$: each has a finite
          weight $w_i > 0$.  The solver maximises
          $\sum_{i : c_i \text{ satisfied}} w_i$.
\end{itemize}

For HCORAP (Section~\ref{sec:sat}), the mapping is direct:
\begin{itemize}
    \item $\varphi_H$ = all hard clauses (reification, AMO, coverage,
          max-hours).
    \item Soft clauses = expertise reward (O1) + continuity penalty (O2) +
          extra-hour penalty (O3), with respective weights.
\end{itemize}

The MaxSAT solver finds an assignment $\sigma^*$ that satisfies all hard
clauses and maximises the total weight of satisfied soft clauses in a
\textbf{single invocation}---no manual binary/linear search is needed.

\subsubsection{Internal Algorithms of MaxSAT Solvers}

State-of-the-art MaxSAT solvers use one of the following internal strategies:

\paragraph{Core-Guided (OLL / RC2).}
The solver iteratively calls a SAT oracle.  When an UNSAT result is obtained,
the solver extracts an \emph{unsatisfiable core}---a minimal subset of soft
clauses that cannot all be satisfied simultaneously.  \emph{Relaxation
variables} are added to the core clauses, and a new cardinality constraint
bounds how many of them can be relaxed.  This process repeats until a SAT
result is found~\cite{rc2maxsat,morgado2013iterative}.

The core-guided approach is particularly effective when the optimum cost (total
weight of violated soft clauses) is small relative to the total weight, as few
cores need to be extracted.

\paragraph{Linear SAT-UNSAT Search.}
Internally equivalent to the top-down linear search
(Algorithm~\ref{alg:topdown}), but executed within the solver with full
learned-clause reuse.

\paragraph{Stratified / Weight-Aware.}
For weighted instances, some solvers partition soft clauses by weight and
solve them in decreasing weight order (stratification), which naturally
handles the lexicographic structure of the HCORAP stratified variant
(Equation~\ref{eq:cov-soft}).

\subsubsection{Stratified Optimisation in HCORAP}

The HCORAP stratified encoding
(Section~\ref{sec:sat}) assigns coverage soft clauses a weight
$W_{\text{total}} + 1$ (where $W_{\text{total}}$ is the sum of all other
soft-clause weights).  This guarantees that the MaxSAT solver first maximises
service coverage, and only then optimises the secondary objective (expertise
reward $-$ continuity penalty $-$ extra-hour cost).  This is equivalent to a
\textbf{two-level lexicographic} MaxSAT problem, which weight-aware solvers
handle natively.

\subsubsection{DIMACS WCNF Format}

The HCORAP tool \texttt{hcorap2sat} outputs instances in the
\textbf{weighted DIMACS CNF (WCNF)} format~\cite{maxsat-eval}, which is the
standard input format for MaxSAT solvers:

\begin{verbatim}
p wcnf <#vars> <#clauses> <top_weight>
<top_weight> <lit_1> <lit_2> ... 0       (hard clause)
<weight>     <lit_1> <lit_2> ... 0       (soft clause)
\end{verbatim}

Any off-the-shelf MaxSAT solver (e.g.\ WMaxCDCL, RC2, CASHWMaxSAT,
EvalMaxSAT, Open-WBO) can then be used to solve the instance.

\subsubsection{Key Properties}

\begin{center}
\begin{tabular}{lc}
\toprule
\textbf{Property} & \textbf{MaxSAT} \\
\midrule
Explicit bound iteration by user  & No (solver-internal) \\
Objective encoding                & Soft clauses with weights \\
Learned clauses reused            & Yes (internally) \\
Optimality guarantee              & Yes (solver proves optimal) \\
Handles weighted objectives       & Yes (natively) \\
Handles stratified / lexicographic & Yes (via weight hierarchy) \\
Implementation complexity         & Low (single solver call) \\
\bottomrule
\end{tabular}
\end{center}


%-----------------------------------------------------------------------------
\subsection{Comparison of Approaches}\label{sec:opt-comparison}
%-----------------------------------------------------------------------------

\begin{table}[ht]
\centering
\caption{Comparison of optimisation approaches for HCORAP.}
\label{tab:opt-comparison}
\small
\begin{tabular}{@{}lccc@{}}
\toprule
\textbf{Criterion}
    & \textbf{Non-Incremental SAT}
    & \textbf{Incremental SAT}
    & \textbf{MaxSAT} \\
\midrule
Number of SAT calls
    & $O(\log W)$ (binary)
    & $O(\log W)$ (binary)
    & 1 (internal iteration) \\[3pt]
Learned-clause reuse
    & \textsf{No}
    & \textsf{Yes}
    & \textsf{Yes} \\[3pt]
Encoding rebuilt?
    & Each call
    & Once
    & Once \\[3pt]
Multi-component weighted obj.
    & Manual (encode sum)
    & Manual (encode sum)
    & Native (soft weights) \\[3pt]
Stratified / lexicographic
    & Manual (two-phase)
    & Manual (two-phase)
    & Native (weight hierarchy) \\[3pt]
Solver infrastructure
    & Any SAT solver
    & IPASIR-compatible
    & MaxSAT solver \\[3pt]
Implementation effort
    & Low
    & Medium
    & Low \\[3pt]
Optimality guarantee
    & Yes (exhaustive search)
    & Yes (exhaustive search)
    & Yes (solver-certified) \\
\bottomrule
\end{tabular}
\end{table}

\paragraph{Discussion.}
For HCORAP specifically, the \textbf{MaxSAT approach} is the most natural fit
because:
\begin{enumerate}[label=(\alph*)]
    \item The problem already has a multi-component weighted objective that maps
          directly to weighted soft clauses.
    \item The stratified variant (lexicographic coverage $\succ$ revenue) is
          handled natively by weight-aware MaxSAT solvers.
    \item No manual search loop is needed; the solver handles all internal
          iteration, core extraction, and optimality certification.
    \item The HCORAP tool can output standard WCNF files, allowing any
          state-of-the-art MaxSAT solver to be used as a black box.
\end{enumerate}

\noindent
The \textbf{incremental SAT approach} (assumption-based binary search) is a
strong alternative when a MaxSAT solver is not available or when finer control
over the search process is desired (e.g.\ for anytime behaviour or custom
heuristics).

\noindent
The \textbf{non-incremental SAT approach} provides the simplest baseline
implementation---useful for prototyping and for instances small enough that the
overhead of rebuilding the encoding is negligible.


%=============================================================================
\section{Dataset}\label{sec:dataset}
%=============================================================================

\subsection{Overview}

The benchmark dataset used in~\cite{collcaballero2025hcorap} consists of
\textbf{400 randomly generated instances} modelling realistic home care
scenarios in the Barcelona metropolitan area.

Each instance is characterised by three parameters:
\begin{itemize}
    \item $U \in \{30, 40\}$: number of users,
    \item $A \in \{10, 15, 20, 25\}$: number of agents,
    \item $V \in \{4, 5\}$: number of services per user (so $S = U \times V$).
\end{itemize}
This gives $2 \times 4 \times 2 = 16$ parameter configurations.
For each configuration, \textbf{50 instances} are generated, yielding
$16 \times 50 = 800$ total in two groups (U30 and U40), with \textbf{400
instances per user count}.

\subsection{Instance Generation}

The instances are generated by the Python script
\texttt{create\_instances.py}~\cite{hcorap-github}, which simulates realistic
home care data:

\begin{itemize}
    \item \textbf{Agents} are characterised by age (25--60), qualifications
          (drawn from \{basic, nurse, CPR, physio, doctor\}), gender, languages,
          race, and geographic location (GPS coordinates from Barcelona).
          Working schedules ($\TSA$) are generated as contiguous blocks of
          4--12~hours per day, repeated over 5~working days.

    \item \textbf{Users} are elderly individuals (age 70--90) with demographic
          attributes (gender, language, race, location).

    \item \textbf{Services} are assigned to users ($V$ services per user) with
          types drawn from \{basic (more frequent), nurse, CPR, physio, doctor\}.
          Each service has a time window of 2--5~consecutive time slots on one
          day of the week.

    \item \textbf{Reward matrix} ($r$): a similarity score between each
          agent--service pair is computed based on qualification match, shared
          demographics (gender, race, language), and geographic distance.
          The raw scores are then discretised into quartiles, yielding values in
          $\{1, 2, 3, 4\}$ for valid assignments and $0$ for forbidden ones
          (when the agent lacks the required qualification).

    \item \textbf{Continuity-of-care groups} ($\SEQ$): services for the same user
          with the same type are grouped into sequences.

    \item \textbf{Working hours:} $\HN_a = \min(\text{total available slots}, 35)$
          and $\HE_a = \max(\text{total available slots} - 35, 0)$.
          The penalty $P = -1$.

    \item \textbf{Time slots:} $T = 12 \times 5 = 60$ (12~one-hour slots per day,
          5~days per week).
\end{itemize}

\subsection{Instance File Format}

Instances are stored in a plain-text format with the following sections:

\begin{center}
\small
\begin{tabular}{ll}
\toprule
\textbf{Section} & \textbf{Content} \\
\midrule
\texttt{\#U}       & Number of users $U$ \\
\texttt{\#S}       & Number of services $S$ \\
\texttt{\#A}       & Number of agents $A$ \\
\texttt{\#TS}      & Number of time slots $T$ \\
\texttt{\#SU}      & Service groups per user (one line per user) \\
\texttt{\#SEQ}     & Continuity-of-care sequences (one line per group) \\
\texttt{\#TSA(i)}  & Agent availability matrix ($A \times T$) \\
\texttt{\#TSS(i)}  & Service time-slot suitability matrix ($S \times T$) \\
\texttt{\#r(i,j)}  & Reward matrix ($A \times S$) \\
\texttt{\#P}       & Extra-hour penalty \\
\texttt{\#HN(i)}   & Normal working hours per agent \\
\texttt{\#HE(i)}   & Extra working hours per agent \\
\bottomrule
\end{tabular}
\end{center}

\subsection{Instance Statistics}

\begin{table}[ht]
\centering
\caption{Summary of instance configurations.}
\label{tab:instances}
\begin{tabular}{cccccc}
\toprule
$U$ & $A$ & $V$ & $S = U \times V$ & $T$ & \# Instances \\
\midrule
30 & 10 & 4 & 120 & 60 & 50 \\
30 & 10 & 5 & 150 & 60 & 50 \\
30 & 15 & 4 & 120 & 60 & 50 \\
30 & 15 & 5 & 150 & 60 & 50 \\
30 & 20 & 4 & 120 & 60 & 50 \\
30 & 20 & 5 & 150 & 60 & 50 \\
30 & 25 & 4 & 120 & 60 & 50 \\
30 & 25 & 5 & 150 & 60 & 50 \\
\midrule
40 & 10 & 4 & 160 & 60 & 50 \\
40 & 10 & 5 & 200 & 60 & 50 \\
40 & 15 & 4 & 160 & 60 & 50 \\
40 & 15 & 5 & 200 & 60 & 50 \\
40 & 20 & 4 & 160 & 60 & 50 \\
40 & 20 & 5 & 200 & 60 & 50 \\
40 & 25 & 4 & 160 & 60 & 50 \\
40 & 25 & 5 & 200 & 60 & 50 \\
\midrule
\multicolumn{5}{l}{\textbf{Total}} & \textbf{800} \\
\bottomrule
\end{tabular}
\end{table}


%=============================================================================
\section*{Acknowledgements}
%=============================================================================

The HCORAP benchmark dataset and MaxSAT encoding are provided by Jordi Coll
Caballero et al.\ and are publicly available on GitHub~\cite{hcorap-github}.


%=============================================================================
% Bibliography
%=============================================================================

\begin{thebibliography}{99}

\bibitem{collcaballero2025hcorap}
Jordi Coll, Jes\'us Gir\'aldez-Cru, Jordi Levy, Zara Calder\'on, Daniel
Gara\~nena, Simone Mancini, and Pedro Meseguer.
\newblock Optimizing Resource Allocation in Home Care Services using MaxSAT.
\newblock In \emph{Proceedings of the International Conference on the
Integration of Constraint Programming, Artificial Intelligence, and Operations
Research (CPAIOR)}, 2025.

\bibitem{hcorap-github}
Jordi Coll Caballero.
\newblock HCORAP -- Home Care Optimised Resource Allocation Problem.
\newblock \url{https://github.com/jordicollcaballero/HCORAP}, 2024.

\bibitem{maxsat-eval}
MaxSAT Evaluations.
\newblock \url{https://maxsat-evaluations.github.io/2022/rules.html}, 2022.

\bibitem{sinz2005cardinality}
Carsten Sinz.
\newblock Towards an Optimal CNF Encoding of Boolean Cardinality Constraints.
\newblock In \emph{Proceedings of the 11th International Conference on
Principles and Practice of Constraint Programming (CP)},
pp.~827--831, Springer, 2005.

\bibitem{bailleux2003totalizer}
Olivier Bailleux and Yacine Boufkhad.
\newblock Efficient CNF Encoding of Boolean Cardinality Constraints.
\newblock In \emph{Proceedings of the 9th International Conference on
Principles and Practice of Constraint Programming (CP)},
pp.~108--122, Springer, 2003.

\bibitem{asin2011cardinality}
Roberto As\'in, Robert Nieuwenhuis, Albert Oliveras, and Enric
Rodr\'iguez-Carbonell.
\newblock Cardinality Networks: A Theoretical and Empirical Study.
\newblock \emph{Constraints}, 16(2):195--221, 2011.

\bibitem{ipasir}
Tom\'a\v{s} Balyo, Armin Biere, Markus Iser, and Carsten Sinz.
\newblock SAT Race 2015.
\newblock \emph{Artificial Intelligence}, 241:45--65, 2016.
\newblock (Defines the IPASIR incremental SAT solver interface.)

\bibitem{rc2maxsat}
Alexey Ignatiev, Antonio Morgado, and Joao Marques-Silva.
\newblock RC2: An Efficient MaxSAT Solver.
\newblock \emph{Journal on Satisfiability, Boolean Modeling and Computation},
11(1):53--64, 2019.

\bibitem{morgado2013iterative}
Antonio Morgado, Carmine Dodaro, and Joao Marques-Silva.
\newblock Core-Guided MaxSAT with Soft Cardinality Constraints.
\newblock In \emph{Proceedings of the 20th International Conference on
Principles and Practice of Constraint Programming (CP)},
pp.~564--573, Springer, 2014.

\end{thebibliography}

\end{document}
